\section{Related work}
\label{sec:related}

\paragraph{Degrees of autonomy, and where knowledge enters.} Existing systems can
be ordered by how much of the modelling decision they leave to the model:
hyperparameter optimisers, where language models mainly address the cold start;
pipeline composers searching jointly over features, model and hyperparameters;
and adaptive planners that revise the search plan as validation feedback arrives.
Orthogonally they differ in where knowledge of libraries and methods comes from.
Some rely on the parametric memory of the backbone; some retrieve auxiliary
material, for instance strategic insights mined from competition
write-ups~\citep{kulibaba2025kompeteai}; others embed API-level knowledge in the
decision process, as MLZero does by combining semantic and episodic memory to
suppress hallucinated and outdated API calls~\citep{fang2025mlzero}.
Section~\ref{sec:axis} argues that this second axis conflates two interventions
with different mechanisms.

\paragraph{Ablations that disagree.} Multi-agent frameworks for competition-style
data science are
numerous~\citep{li2024autokaggle,trirat2025automlagent,kulibaba2025kompeteai,fang2025mlzero},
and their reported ablations disagree instructively. On MLE-bench, search
structure is the largest single contributor in one system and has no measurable
effect in another, with memory ranked last in one study and level with search in
a second~\citep{toledo2025aira}. That is expected when ablations are run one
factor at a time on different backbones and task sets, and the factors are also
shown to interact: advanced search policies gain nothing when restricted to a
weaker operator set~\citep{toledo2025aira}.

\paragraph{Evaluation against tuned baselines.} Agentic AutoML papers that
compare against classical methods almost always compare against defaults. The
closest work in domain is GRACE-DS~\citep{tsymbalov2026graceds}, which evaluates
LLM-powered AutoML agents on tabular tasks under a structured environment and
reports a ``capability equalizer'' effect: structuring the interaction compresses
the spread across backbones. Two differences matter. Its normalised scale is
anchored on an \emph{untuned} HistGradientBoosting oracle, so its unit of
``solved'' sits well below the tuned references we use; and its eight backbones
are frontier and mixture-of-experts models in the 80B--1.6T range, with no dense
small model evaluated. Its result and ours are two sides of one observation
rather than the same claim: GRACE-DS varies the backbone under one structured
environment and finds the spread between backbones compressed, while we vary the
structure under one backbone and find the spread between architectures wide.
Section~\ref{sec:modelaxis} gives the sharper form of the first half---under the
multi-agent system a stronger backbone moves the normalised score by $0.04$,
against the $0.49$ that separates the architectures at a fixed backbone.

\paragraph{Scaffolding against model capability.} Whether scaffolding can
substitute for capability has been asked directly in the embodied setting, where
reasoning, memory, reflection and reinforcement modules were crossed over a
capability ladder of dense models, concluding that external scaffolding can
compensate for weaker native long-horizon
reasoning~\citep{agentspec2026}. We take that design as precedent and ask the
same question where a strong non-agentic reference exists, which supervised
tabular learning provides and embodied benchmarks do not. Agentic coding has
produced parallel evidence that the harness is a large hidden variable: across
configurable harnesses evaluated identically, aggregate score spans $23.8$
points~\citep{harnessbench2026}, and harness choice induces up to a fortyfold
difference in tokens per solved task while the model-to-model pass-rate spread
stays within eight points~\citep{scaffoldeffect2026}. That literature varies the
harness with the task family fixed; we vary the architectural weight of the
system with the family fixed \emph{and} a tuned non-agentic reference available
on the same split, which turns a spread into a calibrated fraction of an
achievable range. A third result bears on our design: when a coding harness is
evolved automatically, ablations localise the gain to tools, middleware and
long-term memory rather than to the system prompt~\citep{harness2026prompt}. Our
prescribed-knowledge axis is a prompt-level intervention, which is why
Section~\ref{sec:axis} treats it as an object of measurement rather than a design
pattern, and our finding that its two halves are sub-additive is what that
localisation predicts.
